\section{Related Work}


\paragraph{Consistency Models.} Consistency models~\citep{cm, ict} are a new family of generative models designed for efficient generation with few model steps without the need for adversarial training. CMs do not rely on a pretrained diffusion model (DM) to generate training targets but instead leverage an unbiased score estimator.
%
Training CMs have been extended to multi-step sampling~\citep{kim2024consistency,wang2024phased,heek2024multistep}, latent space models~\citep{luo2023latent}, ControlNet~\citep{xiao2023ccm}, video~\citep{wang2024animatelcm}, and combined with an additional adversarial loss~\citep{kim2024consistency,kong2023act}.
%
Despite their sampling efficiency, CMs are typically more challenging to train and require significantly more compute resources than their diffusion counterparts. Our work substantially improves the training efficiency of CMs, reducing the cost of future research and deployment on CMs.
% 
Parallel to our work, \citet{heek2024multistep,wang2024phased} relax the single-step constraint of CMs and generalize them to multi-step variants, drawing inspiration from TRACT~\citep{berthelot2023tract}.
Under the piecewise linear design, increasing the number of segments will also lead to diffusion pretraining, showing a different approach to connect CMs to the standard diffusion models.

\paragraph{Diffusion Distillation.} 
%
Drawing inspiration from knowledge distillation~\citep{hinton2015distilling}, distillation is the most widespread training-based approach to accelerate the diffusion sampling procedure. In diffusion distillation, a pretrained diffusion model (DM), which requires hundreds to thousands of model evaluations to generate samples, acts as a teacher. A student model is trained to match the teacher model's sample quality, enabling it to generate high-quality samples in a few steps.

There are two main lines of work in this area. The first category involves trajectory matching, where the student learns to match points on the teacher's sampling trajectory. Methods in this category include 
offline distillation~\citep{Luhman2021KnowledgeDI,dfno,geng2024one}, which require an offline synthetic dataset generated by sampling from a pretrained DM to distill a teacher model into a few-step student model; 
progressive distillation~\citep{salimans2022progressive,meng2023distillation}, and TRACT~\citep{berthelot2023tract}, which require multiple training passes or offline datasets to achieve the same goal; 
and BOOT~\citep{gu2023boot}, Consistency Distillation (CD)\citep{cm}, and Imagine-Flash~\citep{kohler2024imagine}, which minimize the difference between the student predictions at carefully selected points on the sampling trajectory.

CD is closely related to our method, as it leverages a teacher model to generate pairs of adjacent points and enforces the student predictions at these points to map to the initial data point. However, it employs a fixed schedule derived from a specific sampler, which may introduce non-negligible discretization errors in approximating the consistency condition. It also limits the quality of consistency models to that of the pretrained diffusion model. 
% The LPIPS-based metric~\citep{lpips} used in CD loss introduces undesirable bias when evaluated on metrics such as FID.

The second category minimizes the probabilistic divergence between data and model distributions, \ie distribution matching~\citep{poole2022dreamfusion,wang2024prolificdreamer,di,yin2023one,sid}. These methods~\citep{di,yin2023one,nguyen2023swiftbrush,sauer2023adversarial,kohler2024imagine,xu2023ufogen,lin2024sdxl,zhou2024long} use score distillation or adversarial loss, to distill an expensive teacher model into an efficient student model. However, they can be challenging to train in a stable manner due to the alternating updating schemes from either adversarial or score distillation. Some of these methods such as DreamFusion~\citep{poole2022dreamfusion} and ProlificDreamer~\citep{wang2024prolificdreamer} are used for 3D object generation.  

A drawback of training-based approaches is that they need additional training procedures after pretraining to distill an efficient student, which can be computationally intensive. For a detailed discussion on the recent progress of diffusion distillation, we direct the readers to \citet{dieleman2024distillation}.

\vspace{-5pt}
\paragraph{Fast Samplers for Diffusion Models.}
Fast samplers are usually training-free and use advanced solvers to simulate the diffusion stochastic differential equation (SDE) or ordinary differential equation (ODE) to reduce the number of sampling steps. These methods reduce the discretization error during sampling by analytically solving a part of SDE or ODE~\citep{lu2022dpm,lu2022dpmplus,xue2024sa}, by using exponential integrators and higher order polynomials for better approximation of the solution~\citep{zhang2022fast}, using higher order numerical methods~\citep{Karras2022edm}, using better approximation of noise levels during sampling~\citep{kong2021fast}, correcting predictions at each step of sampling~\citep{zhao2024unipc} and ensuring that the solution of the ODE lies on a desired manifold~\citep{liu2022pseudo}. 
Another orthogonal strategy is the parallel sampling process~\citep{pokle2022deep,shih2024parallel}, solving fixed points of the entire ODE/SDE trajectories. A drawback of these fast samplers is that the quality of samples drastically reduces as the number of sampling steps falls below a threshold such as 10 steps.
